% ===================================================================================
% Appendix C (chunk C): the a priori bound at the hard edge, and Q in I + dC(L^2).
% Labels: app:apriori, prop:apriori, cor:L2, lem:sumbound, lem:lenssector.
% ===================================================================================
\providecommand{\eps}{\varepsilon}
\providecommand{\re}{\operatorname{Re}}
\providecommand{\im}{\operatorname{Im}}
\providecommand{\Res}{\operatorname*{Res}}
\providecommand{\thlip}{\vartheta_{\mathrm{lip}}}
\providecommand{\Ob}{\Omega_{\mathrm{b}}}
\providecommand{\Osat}{\Omega^{\triangle}}
\providecommand{\Oinf}{\Omega_{\infty}}
\providecommand{\Cauchy}{\mathcal{C}}
\providecommand{\dC}{\partial\mathcal{C}}

\section{The a priori bound at the hard edge}\label{app:apriori}
% src: RH_px §5 "(Y4)_x" (one-sided nodes; review REVIEW_RHpx F9); ROUTE2 §4.3 "(Q in L^2)" (Potter loss);
% RH_plt1b §2 (lens sector needs only phi_0 < pi p); ROUTE2_xframe1 §2 row "§5 (Y4)_x"; RH_zeta_close §1.6 (vertex list).

The comparison argument of \S\ref{sec:hardedge} applies Theorem~4.4 of Conway and Deift (\ref{tool:T3}) to $Q$. Its hypothesis is
$Q_\pm\in I+\dC(L^2(\Sigma))$: the boundary values of $Q-I$ must be Cauchy integrals of $L^2$ functions. We never solve a
Riemann--Hilbert problem for $Q$; $Q$ is built from the orthogonal polynomials, and we only have to check that it is not too
singular at the vertices of $\Sigma$. The only delicate vertex is the hard edge $0$, where infinitely many nodes accumulate.
There we bound the Cauchy sums in $Y$ directly (Lemma~\ref{lem:sumbound}), except in thin sectors around the band, where the
poles are removed by $R$ and a contour argument is needed (Lemma~\ref{lem:lenssector}). Throughout this appendix $n\ge n_0$ is
\emph{fixed}, and constants may depend on $n$.

Recall from Lemma~\ref{lem:interp} that
\begin{equation}\label{eq:U-sum}
  U_{j2}(z)=\sum_{k\ge1}\frac{w(X_k)\,U_{j1}(X_k)}{z-X_k}\qquad(j=1,2),\qquad w(X_k)=n^pw_k,
\end{equation}
and that $U_{11},U_{21}$ are polynomials. Put $\varphi_1:=\frac12\thlip$.

\begin{lemma}[Sum bound]\label{lem:sumbound}
For every $\eta>0$ there is $C$ (depending on $n$ and $\eta$) such that $|U_{j2}(z)|\le C|z|^{-\eta}(1+|z|^{\alpha})$ for
$0<|z|\le1$, $\varphi_1\le|\arg z|\le\pi$.
\end{lemma}

\begin{proof}
For such $z$ and $X>0$, $|z-X|\ge\sin\varphi_1\max(|z|,X)$. By Lemma~\ref{lem:potter} (or Lemma~\ref{lem:pottersector}(i)),
$w_k\le C_\eta k^{-s+\eta}$, and $U_{j1}$ is bounded on $[0,n^p]$. Let $K:=n|z|^{-1/p}$, so that $X_k\ge|z|$ if and only if
$k\le K$. Then
\[
  |U_{j2}(z)|\le C\sum_{k\le K}\frac{k^{-s+\eta}}{(n/k)^p}+\frac C{|z|}\sum_{k>K}k^{-s+\eta}
  \le C\bigl(1+K^{1+p-s+\eta}+\log K\bigr)+C|z|^{-1}K^{1-s+\eta}.
\]
Since $K^{1+p-s}=n^{1+p-s}|z|^{\alpha}$ and $|z|^{-1}K^{1-s}=n^{1-s}|z|^{\alpha}$ (recall $(s-1)/p-1=\alpha$), and
$K^\eta=n^\eta|z|^{-\eta/p}$, the claim follows after renaming $\eta$.
\end{proof}

\begin{lemma}[Lens sectors]\label{lem:lenssector}
For every $\eta>0$ there are $b\in(0,\delta)$ and $C$ such that $|R_{j2}(z)|\le C|z|^{-\eta}(1+|z|^\alpha)$ for $0<|z|\le b/2$,
$0<|\arg z|<\varphi_1$, and also for $z\in\Ob\cap U(0,b/2)$.
\end{lemma}

\begin{proof}
Take $z$ with $0<\arg z<\varphi_1$; the case $\arg z<0$ is the mirror image. In $\Ob^+$, $R_{j2}=U_{j2}-U_{j1}We^{-2i\theta}\kappa_+$
by~\eqref{eq:R}. Put
\[
  \mathcal K_j(s):=\frac{U_{j1}(s)W(s)}{e^{2i\theta(s)}-1},\qquad s\in G:=\{s:\ r<|s|<b,\ |\arg s|<\thlip\},
\]
where $b$ and $r$ are chosen with $t(b),t(r)\in\N+\frac12$, and $r<\frac12|z|$, $|z|<\frac12b$. Note $e^{-2i\theta}\kappa_+=1/(e^{2i\theta}-1)$,
so $R_{j2}(z)=U_{j2}(z)-\mathcal K_j(z)$. The function $\mathcal K_j$ is meromorphic on a neighbourhood of $\overline G$: $W$ is
analytic there (Lemma~\ref{lem:weightdomain}), and $e^{2i\theta}=1$ only where $t\in\mathbb Z$, i.e.\ at the nodes (the points with
$t\in-\N$ have $|\arg s|=\pi p>\thlip$). At a node, $e^{2i\theta(s)}-1=2i\theta'(X_k)(s-X_k)(1+O(s-X_k))$, so
$\Res_{X_k}\mathcal K_j=U_{j1}(X_k)W(X_k)/(2i\theta'(X_k))=w(X_k)U_{j1}(X_k)$. The residue theorem on $\partial G$ (positively
oriented) and~\eqref{eq:U-sum} give
\begin{equation}\label{eq:R-contour}
  R_{j2}(z)=\sum_{X_k\notin G}\frac{w(X_k)U_{j1}(X_k)}{z-X_k}-\frac1{2\pi i}\oint_{\partial G}\frac{\mathcal K_j(s)\diff s}{s-z}.
\end{equation}
We bound the terms. (1) The nodes with $X_k\ge b$ are finitely many and at distance $\ge b/2$ from $z$; those with $X_k\le r$
contribute at most $\frac2{|z|}\sum_{X_k\le r}w(X_k)|U_{j1}(X_k)|$, which tends to $0$ as $r\to0$. (2) On the ray
$\arg s=\thlip$, $|e^{2i\theta}-1|^{-1}=|e^{-2i\theta}|/|1-e^{-2i\theta}|\le2\exp(-2\pi n|s|^{-1/p}\sin(\thlip/p))$ by~\eqref{eq:imtheta},
which beats $|W(s)|\le C|s|^{\alpha-\eta}$ (Lemmas~\ref{lem:bandweight}, \ref{lem:pottersector}(ii)); and $|s-z|\ge|s|\sin(\thlip-\varphi_1)$.
So this part is bounded. (3) On the ray $\arg s=-\thlip$ we write $1/(e^{2i\theta}-1)=-1-e^{2i\theta}/(1-e^{2i\theta})$, where now
$|e^{2i\theta}|<1$ is super-small, so $\mathcal K_j=-U_{j1}W+(\text{bounded})$. Since $|s-z|\ge c'(|s|+|z|)$,
\[
  \int_0^b\frac{|W(\rho e^{-i\thlip})|\diff\rho}{\rho+|z|}\le C\int_0^b\frac{\rho^{\alpha-\eta}\diff\rho}{\rho+|z|}\le C|z|^{-\eta}(1+|z|^{\alpha}),
\]
because $\alpha-\eta>-1$. (4) On the arcs $|s|=r$ and $|s|=b$, $|t|\in\N+\frac12$. For such $t=u+iv$ we have
$|e^{2\pi it}-1|\ge c'>0$: if $|v|\ge\frac14$ this is clear, and if $|v|<\frac14$ then $u$ is at distance at least $\frac38$ from
$\mathbb Z$, so $|e^{2\pi it}-1|=2e^{-\pi v}|\sin\pi t|\ge2e^{-\pi/4}\sin(3\pi/8)$. Hence $|\mathcal K_j|\le C|W|\le C|s|^{\alpha-\eta}$
there. The arc $|s|=r$ has length $\le2\thlip r$ and $|s-z|\ge\frac12|z|$, so its contribution is $\le Cr^{1+\alpha-\eta}/|z|\to0$
as $r\to0$; the arc $|s|=b$ contributes a bounded amount. Letting $r\to0$ in~\eqref{eq:R-contour} proves the bound for
$0<\arg z<\varphi_1$. For $z\in\Ob\cap U(0,b/2)$ with $\varphi_1\le|\arg z|<\frac54\thlip$, $R_{j2}=U_{j2}-U_{j1}We^{\mp2i\theta}\kappa_\pm$
is bounded by Lemma~\ref{lem:sumbound} plus a super-small term.
\end{proof}

\begin{proposition}[A priori bound]\label{prop:apriori}
Fix $n\ge n_0$. Then $\det Q\equiv1$, and:
\begin{enumerate}
\item[(i)] for every $\eta>0$ there is $C_{n,\eta}$ with $|Q(z)|\le C_{n,\eta}|z|^{-|\alpha|/2-\eta}$ for $z\in U(0,\delta)\setminus\Sigma$;
\item[(ii)] $Q$ is bounded on every bounded set at positive distance from $0$, and $Q(z)\to I$ as $z\to\infty$;
\item[(iii)] $Q$ has continuous boundary values on $\Sigma$ minus its vertices.
\end{enumerate}
The same holds for $Q^{-1}$.
\end{proposition}

\begin{proof}
$\det U=\det Y=1$ (Lemma~\ref{lem:interp}), and every other factor in~\eqref{eq:VM} and in $Q=D_\infty^{-\sigma_3}VD_W^{\sigma_3}$ has
determinant $1$; so $\det Q=1$ and $Q^{-1}$ has the same entries up to sign and order.

(i) Let $\eta>0$ and $|z|<\min(\delta,b/2)$ (the annulus $b/2\le|z|<\delta$ is covered by~(ii)). Write
$Q=E_\infty\,U\,M\,D_W^{\sigma_3}$ with the constant $E_\infty:=D_\infty^{-\sigma_3}C_\infty e^{-n\gamma_p\sigma_3}$ and $M$
from~\eqref{eq:VM}. We use: $|e^{\pm n\varphi}|\le e^{Cn|z|^{1/2}}\le C$ (Lemma~\ref{lem:signs}(iii)); $|F|^{\pm1}\le C$ outside
the lens and $|\tilde F|^{\pm1}\le C$ in the lens near $0$ (O5); $|D_W(z)|^{\pm1}\le C|z|^{\pm\alpha/2-\eta}$
(Lemma~\ref{lem:szego}(vi)); $|1/W(z)|\le C|z|^{-\alpha-\eta}$ in the lens (Lemmas~\ref{lem:bandweight},
\ref{lem:pottersector}(ii)); and $U_{j1}$ bounded.

\emph{Outside the lens} ($z\in\Oinf$), $|\arg z|\ge\frac34\thlip>\varphi_1$ and $M=e^{n\varphi\sigma_3}F^{\sigma_3}$. The first column of
$Q$ is a constant times $U_{\cdot1}e^{n\varphi}FD_W$, of size $\le C|z|^{\alpha/2-\eta}$; the second is a constant times
$U_{\cdot2}e^{-n\varphi}F^{-1}D_W^{-1}$, of size $\le C|z|^{-\eta}(1+|z|^\alpha)|z|^{-\alpha/2-\eta}$ by Lemma~\ref{lem:sumbound}. Both are
$\le C|z|^{-|\alpha|/2-2\eta}$.

\emph{In the lens} $\Ob^\pm$, $UM=R\,e^{n\varphi\sigma_3}A_\pm\tilde F^{\sigma_3}$, and
$e^{n\varphi\sigma_3}A_\pm=\bigl(\begin{smallmatrix}e^{n\varphi}&0\\\mp e^{n\varphi}/W&e^{-n\varphi}\end{smallmatrix}\bigr)$. The first column of
$R$ is that of $U$, and the second is bounded by Lemma~\ref{lem:lenssector}. Hence the first column of $Q$ is at most
$C(1+|z|^{-\alpha-\eta}|z|^{-\eta}(1+|z|^\alpha))|z|^{\alpha/2-\eta}\le C|z|^{-|\alpha|/2-3\eta}$, and the second is at most
$C|z|^{-\eta}(1+|z|^\alpha)|z|^{-\alpha/2-\eta}\le C|z|^{-|\alpha|/2-2\eta}$. Renaming $\eta$ gives~(i).

(ii) Away from $0$ each of $U$ (meromorphic, with finitely many nodes in any compact set not containing $0$), the triangular
and diagonal factors in~\eqref{eq:R}--\eqref{eq:V}, $F$, $\tilde F$ and $D_W$ is analytic in each region up to its boundary, except
for the poles of $U$ and of $\kappa_\pm$ at the nodes, which cancel in $R$ (Step~1 of \S\ref{sec:frame-transf}). Near $4c$ the
only singular factors are $\tilde F\sim(4c-z)^{-1/2}$ and $D_W\sim(4c-z)^{1/2}$, and they occur only in the product
$\tilde FD_W$, whose square $(\tilde FD_W)^2=C_nz^\alpha r_n\omega$ is bounded above and below on $U(4c,\eps)$
(Lemma~\ref{lem:szego}(vi)(b)). At the remaining vertices ($4c\pm ih_0$, $n^p+\delta'\pm ih_0$, $n^p+\delta'$) all factors are
analytic and bounded. At $\infty$, $Q\to I$ by Lemma~\ref{lem:normalisation}. (iii) follows from the same description and
Lemma~\ref{lem:jumpsV}.
\end{proof}

\begin{corollary}[$Q$ is an admissible solution]\label{cor:L2}
For each fixed $n\ge n_0$, $f:=Q_+-Q_-\in L^2(\Sigma)$, $Q_\pm=I+\Cauchy_\pm f$, and $Q=I+\Cauchy f$ off $\Sigma$. Thus
$Q_\pm\in I+\dC(L^2(\Sigma))$ in the sense of \ref{tool:T3}, and the same holds for $Q^{-1}$.
\end{corollary}

\begin{proof}
For fixed $n$, $\Sigma$ is bounded, $|Q_\pm|$ is bounded away from $0$, and $|Q_\pm(z)|\le C|z|^{-(1+|\alpha|)/4}$ near $0$ by
Proposition~\ref{prop:apriori} with $\eta=(1-|\alpha|)/4$. Since $(1+|\alpha|)/4<\frac12$, $f\in L^2(\Sigma)$. Put
$H:=Q-I-\Cauchy f$ on $\C\setminus\Sigma$.

\emph{$H$ has no jump.} On each open arc of $\Sigma$, $Q_\pm$ are real-analytic (each region's formula for $Q$ continues
analytically across the arc, the poles at the nodes having been removed), so $f$ is locally H\"older and $\Cauchy f$ has
continuous boundary values with $\Cauchy_+f-\Cauchy_-f=f$ (Sokhotski--Plemelj). Hence $H_+=H_-$ on open arcs, and $H$ is
analytic across them by Morera's theorem. So $H$ is analytic in $\C$ minus the finitely many vertices.

\emph{The vertices are removable.} Let $z_0$ be a vertex. By Cauchy--Schwarz and $|\Sigma\cap U(z,\rho)|\le C\rho$,
$|\Cauchy f(z)|\le\|f\|_2(\int_\Sigma|s-z|^{-2}|\diff s|)^{1/2}\le C\|f\|_2\,\mathrm{dist}(z,\Sigma)^{-1/2}$. Near $z_0$, $\Sigma$
consists of finitely many $C^2$ arcs, so $\iint_{\rho<|z-z_0|<2\rho}\mathrm{dist}(z,\Sigma)^{-1/2}\diff A=O(\rho^{3/2})$, while
$\iint_{\rho<|z-z_0|<2\rho}|Q-I|\diff A=O(\rho^{2-(1+|\alpha|)/4})$. Hence $\iint_{\rho<|z-z_0|<2\rho}|H|\diff A=o(\rho)$. Choosing
$r\in(\rho,2\rho)$ with $\oint_{|z-z_0|=r}|H||\diff z|\le\rho^{-1}\iint_{\rho<|z-z_0|<2\rho}|H|\diff A$, the Laurent coefficients
satisfy $|c_{-k}|\le\frac1{2\pi}r^{k-1}\oint_{|z-z_0|=r}|H||\diff z|=o(\rho^{k-1})$, so $c_{-k}=0$ for $k\ge1$.

\emph{Liouville.} $H$ is entire and tends to $0$ at $\infty$ (both $Q-I$ and $\Cauchy f$ are $O(1/z)$, $\Sigma$ being bounded),
so $H\equiv0$. The boundary values give $Q_\pm=I+\Cauchy_\pm f$. The argument for $Q^{-1}$ is identical.
\end{proof}

\noindent This is the content of~\Fi7. For $\eta$ in~\Fi7 any value in $(0,\frac{1-|\alpha|}2)$ makes the boundary values of $Q$
square integrable at $0$; this is the place where $|\alpha|<1$ enters the a priori bound.
